\chapter{Documentation and developer experience}
\label{ch:devx}

\begin{aside}
A library that is easy to use is not the same as a library that
is easy to learn. \maya{}'s architecture is small, but the
surface area --- DSL, builders, widgets, signals, themes ---
takes a while to absorb. This chapter is about \emph{developer
experience}: how to document a \maya{} library, how to write
examples, how to teach the architecture. Part of the book's job
is to teach by example; here we name the techniques.
\end{aside}

\section{Three audiences}

When you write docs for a \maya{}-based project, you are writing
for three audiences:

\begin{itemize}[leftmargin=*]
  \item \textbf{Users.} People who run your app. They need
    the keyboard shortcuts, the troubleshooting, the upgrade
    guide.
  \item \textbf{Integrators.} People who embed your library
    in their app. They need the API reference, the example
    code, the architectural overview.
  \item \textbf{Contributors.} People who modify your source.
    They need the design rationale, the testing setup, the
    build instructions.
\end{itemize}

A single \code{README.md} cannot serve all three. Split them.

\section{The README}

For \maya{} projects, the README answers ``what is this and
why should I care?'' in 60 seconds:

\begin{itemize}[leftmargin=*]
  \item One-line description.
  \item Screenshot or recorded animation.
  \item Install command.
  \item First-run instructions.
  \item Links to deeper docs.
\end{itemize}

That is it. Long READMEs lose readers. Short READMEs invite the
right next click.

\section{API reference}

For each public type or function, document:

\begin{itemize}[leftmargin=*]
  \item One-sentence purpose.
  \item Parameters: types, units, valid ranges.
  \item Return value.
  \item Errors / exceptions.
  \item Example of typical use.
\end{itemize}

Doxygen extracts most of this from comments. The discipline:
write the comment before the code, not after.

\subsection*{Comment style}

\begin{lstlisting}
/// Render an element to ANSI bytes.
///
/// Lays out the element against the current terminal size,
/// diffs against the previous canvas, and writes the
/// minimum-change byte sequence.
///
/// \\param e   the element to render.
/// \\returns  the number of bytes written.
/// \\throws   io_error on terminal write failure.
size_t render(const Element& e);
\end{lstlisting}

Brief, parameter, return, throws. The Doxygen tags are
optional; the structure matters.

\section{Examples}

The fastest way to learn an API is to see it used. For each
major feature, ship at least one example:

\begin{itemize}[leftmargin=*]
  \item A minimal sample that demonstrates the feature in
    isolation.
  \item A realistic sample showing typical composition.
  \item Tests that double as examples.
\end{itemize}

\maya{}'s \code{examples/} directory holds small one-file
programs: counter, todo, file-browser, dashboard. Each
exercises a slice of the library.

\subsection*{Runnable in seconds}

Examples should build and run without setup. A \code{cmake}
sub-target builds each example; \code{make examples} produces
binaries; the README shows the commands.

\section{Architecture documents}

Beyond per-function docs, your project needs a few
architecture-level documents:

\begin{itemize}[leftmargin=*]
  \item \textbf{Overview.} The big picture in one page.
  \item \textbf{Decision records.} Why we chose X over Y.
  \item \textbf{Concepts.} Vocabulary specific to your app.
  \item \textbf{Sequence diagrams.} For complex flows
    (startup, network sync, crash recovery).
\end{itemize}

The book you are reading is an architecture document for
\maya{} itself. For your own apps, a few-page overview is
usually enough.

\subsection*{Decision records (ADRs)}

A lightweight format for design decisions:

\begin{itemize}[leftmargin=*]
  \item Title.
  \item Context: what was the situation?
  \item Decision: what did we choose?
  \item Consequences: what does this mean for callers?
  \item Status: proposed, accepted, deprecated, superseded.
\end{itemize}

A six-paragraph ADR is enough. Keep them in
\code{docs/adr/}; number them sequentially.

\section{Tutorials}

Tutorials are different from reference. They take a learner
from zero to a working application via a guided sequence:

\begin{itemize}[leftmargin=*]
  \item State assumed knowledge upfront.
  \item Build one feature at a time.
  \item Show the diff at each step.
  \item Explain \emph{why}, not just \emph{how}.
\end{itemize}

A good tutorial leaves the learner with a working app and the
confidence to extend it.

\section{Error message quality}

User-facing errors are docs too. A bad error:

\begin{lstlisting}
error: parse failed
\end{lstlisting}

A good error:

\begin{lstlisting}
error: parse failed at config.toml:14:5
  expected: a string for key 'theme'
  found:    integer 42
  hint:     theme = \"dark\"
\end{lstlisting}

Three things distinguish them: location, expected vs found,
a hint. Build the helper once; use it for every parser, schema,
or validator in your app.

\subsection*{Compile errors}

For libraries that lean on templates, compile errors are
docs. \maya{}'s concept names appear in errors; the names are
chosen to be human-readable. Apply the same discipline to your
own templates.

\section{Changelog}

Every release needs a changelog. The format that wins:

\begin{lstlisting}
## [0.5.0] - 2025-03-15

### Added
- Mouse focus on click.
- New `picker` widget with fuzzy matching.

### Changed
- Default theme now uses 16-colour fallback on dumb terminals.

### Fixed
- Crash on resize during animation.
\end{lstlisting}

Sections by kind, dates in ISO format. The first thing users
read after upgrading; the first thing developers grep before
upgrading.

\section{Onboarding contributors}

The \code{CONTRIBUTING.md}:

\begin{itemize}[leftmargin=*]
  \item Build instructions.
  \item Test instructions.
  \item Style guide.
  \item How to file an issue / open a PR.
  \item Areas that need help.
\end{itemize}

Make the first PR easy. A ``good first issue'' label and a
sample issue with hand-holding goes a long way.

\section{Versioning}

Semantic versioning is the convention:

\begin{itemize}[leftmargin=*]
  \item \code{MAJOR}: incompatible API changes.
  \item \code{MINOR}: new functionality, backward-compatible.
  \item \code{PATCH}: bug fixes.
\end{itemize}

Pin dependencies to compatible versions; bump \code{MAJOR} only
when truly necessary.

\section{Deprecation policy}

Things you deprecate should:

\begin{itemize}[leftmargin=*]
  \item Be announced one minor version before removal.
  \item Print a warning when used.
  \item Have a documented replacement.
  \item Be removed only in a major version.
\end{itemize}

The C++ attribute \code{[[deprecated(\"use foo instead\")]]}
generates compile-time warnings.

\section{Building trust}

Documentation is also a trust signal. Users adopt libraries
that:

\begin{itemize}[leftmargin=*]
  \item Are well-tested (visible CI).
  \item Have a changelog.
  \item Respond to issues.
  \item Have other users (a few public projects using it).
\end{itemize}

You cannot fake any of these; build them over time.

\section{Pitfalls}

\begin{warning}
\textbf{Documenting only the happy path.} The error paths are
often where users get stuck. Document failure modes
explicitly.
\end{warning}

\begin{warning}
\textbf{Stale examples.} Examples that compiled in version
0.1 but not in 0.5 are worse than no examples. CI should
build every example on every commit.
\end{warning}

\begin{warning}
\textbf{Magical docs.} ``Just call \code{run}'' is not
documentation. Explain what \code{run} does, what it requires,
what it produces.
\end{warning}

\begin{warning}
\textbf{Code without comments.} For non-trivial logic, a
sentence explaining \emph{why} saves future-you hours. Avoid
restating \emph{what} the code already says.
\end{warning}

\section*{Workshop}

\begin{enumerate}[leftmargin=*]
  \item Take a function from your project. Write a 4-line
    doc comment: purpose, params, return, error.
  \item Write a 6-paragraph ADR for one design choice in your
    project.
  \item Add an example to your repo. Wire it into CI so it
    builds on every commit.
\end{enumerate}

\begin{recap}
A library is its docs plus its code. Write for three audiences:
users, integrators, contributors. Keep READMEs short; document
APIs with purpose, parameters, returns, errors; ship runnable
examples; record architecture decisions; produce a changelog.
The discipline pays back when someone else adopts your library
or when future-you returns to it.
\end{recap}
